\appendix
\section{Ap\'{e}ndices}
\label{sec:appendices}

% --------------------------------------------------------------------------
% Apendice A: ocupaciones extremas en la distribucion de exposicion
% --------------------------------------------------------------------------

\subsection{Ocupaciones en las colas de la distribuci\'{o}n de exposici\'{o}n a LLMs}
\label{app:peru_extremes}

Las dos tablas a continuaci\'{o}n complementan la \cref{subsec:data_descriptives} y dan contenido sustantivo a las dos colas de la distribuci\'{o}n del puntaje $\beta$ de \citet{Eloundou2023_gpt_exposure} para las ocupaciones presentes en la EPEN 2021.

\begin{table}[H]
\centering
\small
\begin{threeparttable}
\caption{Diez ocupaciones m\'{a}s expuestas a LLMs (CIUO-08 4d), Lima Metropolitana}
\label{tab:peru_top10}
\begin{tabular}{c c l}
\toprule
C\'{o}digo CIUO & $\beta$ & Ocupaci\'{o}n \\
\midrule
2432 & 0.806 & Profesionales en relaciones p\'{u}blicas \\
2643 & 0.741 & Traductores, int\'{e}rpretes y otros ling\"{u}istas \\
2514 & 0.706 & Programadores de aplicaciones \\
2641 & 0.689 & Autores y escritores \\
2120 & 0.682 & Matem\'{a}ticos, actuarios y estad\'{i}sticos \\
3312 & 0.662 & Oficiales de cr\'{e}dito y pr\'{e}stamos \\
2431 & 0.661 & Profesionales de publicidad y marketing \\
3512 & 0.656 & T\'{e}cnicos de soporte de usuarios de TIC \\
4323 & 0.654 & Empleados de transporte \\
2513 & 0.638 & Desarrolladores web y multimedia \\
\bottomrule
\end{tabular}
\begin{tablenotes}[flushleft]
\footnotesize
\item \emph{Notas:} Ranking por valor de $\beta$ de \citet{Eloundou2023_gpt_exposure} sobre las ocupaciones presentes en la EPEN 2021.
\end{tablenotes}
\end{threeparttable}
\end{table}

\begin{table}[H]
\centering
\small
\begin{threeparttable}
\caption{Diez ocupaciones menos expuestas a LLMs (CIUO-08 4d), Lima Metropolitana}
\label{tab:peru_bottom10}
\begin{tabular}{c c l}
\toprule
C\'{o}digo CIUO & $\beta$ & Ocupaci\'{o}n \\
\midrule
3421 & 0.000 & Atletas y deportistas \\
9122 & 0.000 & Limpiadores de veh\'{i}culos \\
9312 & 0.000 & Pe\'{o}nes de obras p\'{u}blicas \\
9313 & 0.002 & Pe\'{o}nes de la construcci\'{o}n de edificios \\
7214 & 0.007 & Montadores de estructuras met\'{a}licas \\
8211 & 0.008 & Ensambladores de maquinaria mec\'{a}nica \\
7131 & 0.013 & Pintores y empapeladores \\
7133 & 0.015 & Limpiadores de fachadas y deshollinadores \\
9412 & 0.016 & Ayudantes de cocina \\
7132 & 0.017 & Pintores con pistola y barnizadores \\
\bottomrule
\end{tabular}
\begin{tablenotes}[flushleft]
\footnotesize
\item \emph{Notas:} Ranking inverso por $\beta$. C\'{o}digo CIUO 9313 (pe\'{o}nes de construcci\'{o}n) concentra una masa de empleo expandido sustancial en Lima Metropolitana, ilustrando que la cola de baja exposici\'{o}n no es residual.
\end{tablenotes}
\end{threeparttable}
\end{table}

% --------------------------------------------------------------------------
% Apendice B: heterogeneidad por formalidad en submuestras
% --------------------------------------------------------------------------

\subsection{Heterogeneidad por formalidad estimada en submuestras}
\label{app:peru_subsamples}

La \cref{tab:peru_formality_split} reporta una verificaci\'{o}n alternativa del an\'{a}lisis interactuado de la \cref{tab:peru_formality}: el panel se divide en la mediana ponderada por empleo de $F_o^{2021}$ y la \cref{eq:did_main} se estima por separado en cada submuestra. El patr\'{o}n cualitativo coincide con la regresi\'{o}n interactuada del cuerpo principal y entrega coeficientes estimables sin imponer la restricci\'{o}n de aditividad entre canales.

\begin{table}[H]
\centering
\small
\begin{threeparttable}
\caption{$\hat\tau^{\text{Per\'{u}}}$ en submuestras divididas por formalidad de l\'{i}nea base}
\label{tab:peru_formality_split}
\begin{tabular}{l c c c}
\toprule
                                                & log empleo               & log salario hora       & horas semanales       \\
\midrule
\multicolumn{4}{l}{\emph{Submuestra A. Alta formalidad de l\'{i}nea base ($F_o^{2021} > 0.225$, $N = 1{,}446$)}} \\
$\hat\tau^{\text{Per\'{u}}}$                    & $\mathbf{0.720^{***}}$    & $-0.077$               & $-2.480^{*}$          \\
                                                & $(0.232)$                 & $(0.055)$              & $(1.262)$             \\
\midrule
\multicolumn{4}{l}{\emph{Submuestra B. Baja formalidad de l\'{i}nea base ($F_o^{2021} \leq 0.225$, $N = 632$)}} \\
$\hat\tau^{\text{Per\'{u}}}$                    & $-0.545$                  & $\mathbf{0.254^{***}}$ & $-4.545$              \\
                                                & $(0.339)$                 & $(0.091)$              & $(4.656)$             \\
\bottomrule
\end{tabular}
\begin{tablenotes}[flushleft]
\footnotesize
\item \emph{Notas:} Cada submuestra se estima con la misma especificaci\'{o}n de la \cref{eq:did_main}, ponderada por empleo, con efectos fijos de ocupaci\'{o}n y tiempo y errores est\'{a}ndar agrupados a nivel de ocupaci\'{o}n. La mediana de corte $F_o^{2021} = 0.225$ se calcula sobre 2021 ponderada por empleo. *** $p<0.01$, ** $p<0.05$, * $p<0.10$.
\end{tablenotes}
\end{threeparttable}
\end{table}
